\section{考察}
\label{sec_discussion}
\subsection{定義の観測可能性とラベルの再現性}
第5版のお題文ありは，境界の規則を書き足すことで，同じ判断者の付け直しで $ \kappa $ 0.95 を得ていた．
提案の属性も付け直しで層の $ \kappa $ が0.95以上，規則で導いたお題文ありの $ \kappa $ が0.92〜1.00であり，意図を記録しなくても再現性は下がらなかった．
両者の違いは，再現性そのものではなく，ラベルの根拠が画像に現れるかにある．
第5版では，同じ見た目の字幕の問いに，元の字幕であればお題文なし，投稿者が加えた文であればお題文ありという異なるラベルを求めていた．
提案の属性では，見た目が同じ画像には同じ属性が付くため，判断の揺れは属性ではなく規則の選択として明示される．
実際に，テロップを含めるかという判断は，$ r_1 $ と $ r_2 $ の比較として，ラベルを付け直さずに評価できた．

\subsection{規則の選択がデータに与える影響}
規則の選択によって，お題文ありの割合は7.40\%から16.67\%まで変わった．
お題文ありは，ボケの生成で画像中の文を題として読む処理を適用する対象を決めるため，この差は下流の処理の対象の数に直接影響する．
$ r_{1b} $ は第5版との差が最も小さく，第5版のお題文ありを用いた既存の分析との互換性を保てる．
一方，$ r_{1b} $ では，余白に空欄を含む文と並ぶ状況の文が，問い，指示，題目の句の形でない限りお題文の根拠にならない．
また，重ねた文字の問いを含めるかは，第5版との一致ではほぼ決まらず，台詞の字幕が問いの形のものを題とみなしてよいかという判断に帰着する．
この判断は，画像だけでなく，その画像に付いたボケが問いに答えているかによって確かめられる可能性がある．

\subsection{教師と生徒への影響}
提案の属性への移行の主な障害は，教師の精度である．
Qwen3-VL-32B は，文の形はある程度答えられたが，文字がどの層にあるかを区別できず，層の誤りが規則を通じてお題文ありの誤りに伝わった．
一方，生徒は余白の文字と吹き出しを埋め込みだけで学べており，層の判断は画像の見た目に現れている．
このことから，層は少数の目視ラベルで学習した生徒に任せ，教師には層をまとめた文の形だけを問う構成が考えられる．
文の形は，1,620枚の目視ラベルでは学習が不足し，学習曲線も頭打ちに近かったため，数万枚規模の教師ラベルが必要と考えられる．
書き起こしの文を生徒に加えると精度は上がったが，文字ありの画像の全てに適用するには，最も速い方法でも約2.1日を要する．

\subsection{限界}
本研究には次の限界がある．
第1に，目視ラベルは全て同じ判断者によるものであり，再現性は同じ判断者の付け直しの一致である．
人どうしの一致は測っていない．
第2に，提案の属性の評価に用いた1,620枚は，第5版の評価のために誤りの疑われる画像を多く選んだ集合を含み，重みなしの値は全データより食い違いが多く出る．
第3に，比較の基準のパッケージは約4.5万枚の教師ラベルで学習しており，1,620枚で学習した生徒との比較は学習の量の差を含む．
